% !TeX spellcheck = de_DE

Die Evaluation verfolgt drei zentrale Ziele, welche sich aus den Forschungsfragen ergeben:
\begin{itemize}
	\item \textbf{Z1 - Eignung der Metriken (F1):} Das erste Ziel verfolgt die Prüfung, inwieweit die ausgewählten Lesbarkeitsmetriken, linguistische Oberflächenmetriken und semantische Metriken (vgl. Kapitel \ref{ch:metrikenanalyse}) Unterschiede zwischen weniger und stärker gelungenen LS+-Texten abbilden können.
	\item \textbf{Z2 - Wirkung des Feedbackloops (F2):} Das zweite Ziel beleuchtet die Untersuchung, ob die Metriken sich im Verlauf der Iterationen in Richtung der definierten Zielbereiche entwickeln und ob sich dadurch eine stabile Textqualität erzielen lässt.
	\item \textbf{Z3 - Auswirkung auf Textqualität (F3):} Das dritte Ziel ist die Analyse, wie sich die iterative Überarbeitung eines Textes auf Verständlichkeit, Regelkonformität und inhaltliche Konsistenz auswirkt.
\end{itemize}